% Generated draft: requires mathematical and source review.
\documentclass[11pt]{article}
\usepackage[margin=25mm]{geometry}
\usepackage{fontspec}
\setmainfont{Latin Modern Roman}
\usepackage{amsmath,amssymb,mathtools,unicode-math}
\setmathfont{Latin Modern Math}
\usepackage{enumitem,needspace,xurl,hyperref}
\setlength{\parindent}{0pt}
\setlength{\parskip}{5pt}
\setlength{\emergencystretch}{3em}
\newcommand{\N}{\mathbb{N}}
\newcommand{\Z}{\mathbb{Z}}
\newcommand{\Q}{\mathbb{Q}}
\newcommand{\R}{\mathbb{R}}
\newcommand{\C}{\mathbb{C}}
\newcommand{\F}{\mathbb{F}}
\newcommand{\A}{\mathbb{A}}
\newcommand{\PP}{\mathbb{P}}
\newcommand{\E}{\mathbb{E}}
\newcommand{\eps}{\varepsilon}
\newcommand{\dd}{\,\mathrm{d}}
\begin{document}
\section*{UnsolvedMath ID: 6700052}
\textbf{Source problem code:} catalogue entry
\subsection{Definitions and mathematical statement}
The source calls this Conjecture 18, the Interior Hemi-Spherical Area Inequality. Let $X$ be a compact Riemannian $n$-manifold with boundary and scalar curvature
\[
\operatorname{Sc}_X\geq\operatorname{Sc}_{S^n}=n(n-1),
\]
where the comparison sphere is unit round. It asserts a compact-notation bound, printed as ``co-s+ar''$(X^{-r})\leq2\pi$ for every $r\geq\pi/2$, and explains that no strictly area-decreasing proper map
\[
f:X^{-r}\longrightarrow S^n_+,\qquad \deg f\ne0,
\]
exists when $r>\pi/2$. The exact invariant notation is not defined in the record; the value $2\pi$ is retained as printed, without guessing its normalization.

The source's $r$-interior refers to points at least $r$ from metric incompleteness: in an incomplete manifold a point belongs if its closed $r'$-balls are compact for every $r'<r$. For a compact manifold with boundary this must be interpreted through its interior, or via distance to the boundary, rather than compactness of balls in the already compact space including its boundary. The record does not specify whether $X^{-r}$ uses a strict or non-strict distance cutoff, an induced metric, or a smooth regularized domain.

For a smooth map between Riemannian manifolds of dimension at least two, a customary strictly area-decreasing condition is
\[
\|\Lambda^2D_xf\|<1\qquad\text{at every }x,
\]
meaning strict contraction of every tangent two-dimensional area. A uniform bound by one fixed constant below one is stronger and is not stated. Proper means inverse images of compact sets are compact. Nonzero degree requires a specified orientation or coefficient convention and an appropriate proper-map or relative-boundary setup.

Whether $S^n_+$ denotes an open or closed hemisphere is missing and matters: compact-domain proper maps to a closed hemisphere need boundary/relative-degree conditions, while the open target uses a different proper-degree setting. The no-map formulation has the strict threshold $r>\pi/2$; its endpoint must not be inferred from the unexplained invariant bound at $r=\pi/2$. The source says the situation is unclear even for domains in closed spin manifolds, then ends at ``For instance:'' without the promised example.
\subsection{Short English statement}
At scalar curvature at least the unit sphere's, can a sufficiently deep interior map properly with nonzero degree to a hemisphere while strictly shrinking area? The source forbids this for $r>\pi/2$, with unresolved invariant, hemisphere and degree conventions.
\subsection{Sources}
\begin{itemize}
\item Local UnsolvedMath registry, permanent ID 6700052, exact source code catalogue entry; proposed by Misha Gromov, 2017. Registry status open, classification OPEN-TRIAGE. Source printed tag [?55], occurrence 54, named Conjecture 18.
\item Misha Gromov, \emph{101 Questions, Problems and Conjectures around Scalar Curvature} (2017), PDF page 58, \url{https://www.ihes.fr/~gromov/wp-content/uploads/2018/08/101-problemsOct1-2017.pdf}. Supplied normalized pdf-text extraction accessed July 29, 2026; incomplete author-hosted PDF. No PDF was accessed here.
\item The supplied summary relates this to $r$-interiors and hemispherical comparison and reports no dedicated proof located. It provides no later theorem or reference. The source's explicit uncertainty about domains in closed spin manifolds is retained.
\item Provenance flags current status and statement-reuse terms NEEDS\_REVIEW; the extract stops before an announced example.
\end{itemize}
\end{document}
